\section{Methods}
\label{sec:methods}
\subsection{Representation and reconstruction objective}
For $N$ tensors $\mathcal X_i\in\mathbb R^{m_1\times\cdots\times m_K}$,
choose dimensions $r_k\leq m_k$ and shared bases
$U_k\in\mathbb R^{m_k\times r_k}$ with $U_k^\top U_k=I$.
CSA encodes and reconstructs each sample as
\begin{equation}
\mathcal Y_i=\mathcal X_i\times_1U_1^\top\cdots\times_KU_K^\top,
\qquad \widehat{\mathcal X}_i=\mathcal Y_i\times_1U_1\cdots\times_KU_K,
\label{eq:encoding}
\end{equation}
where $\times_k$ multiplies along mode $k$. The bases minimize
$E=\sum_i\|\mathcal X_i-\widehat{\mathcal X}_i\|_F^2$.
Orthogonal projection gives
$E=E_0-\sum_i\|\mathcal Y_i\|_F^2$, with constant
$E_0=\sum_i\|\mathcal X_i\|_F^2$; minimizing reconstruction error therefore
maximizes retained core energy \citep[Sec.~II-B]{xu2008csa}.

\subsection{Alternating optimization}
Initialize column-orthonormal bases. Fix all modes except $k$, project along
those fixed modes, and unfold each result into
$B_i^{(k)}\in\mathbb R^{m_k\times\prod_{j\ne k}r_j}$.
The core unfolding is $U_k^\top B_i^{(k)}$, so
\begin{equation}
C_k=\sum_iB_i^{(k)}B_i^{(k)\top},\qquad
E(U_k\mid U_{-k})=E_0-\operatorname{tr}(U_k^\top C_kU_k).
\label{eq:conditional}
\end{equation}
Here $U_{-k}$ denotes the fixed bases. The leading $r_k$ eigenvectors of $C_k$ maximize the trace
\citep[Theorem~1]{xu2008csa}. With eigenvalues in descending order,
\begin{equation}
E_{\mathrm{new}}-E_{\mathrm{old}}
=\operatorname{tr}\!\bigl((U_k^{\mathrm{old}})^\top C_kU_k^{\mathrm{old}}\bigr)
-\sum_{j=1}^{r_k}\lambda_j(C_k)\leq0.
\label{eq:nonincrease}
\end{equation}
The old basis remains feasible for the same conditional problem, so its exact
optimum cannot be worse. Cycle through the modes using the latest bases until
stabilization or an iteration limit. These updates are coupled through $C_k$.
At fixed ranks and preprocessing, the nonnegative objective is non-increasing
and converges; this proves neither global optimality nor convergence of the
bases. Appendix~\ref{app:proof} supplies the detailed argument.

\subsection{Course connections and recognition}
With consistent column-major vectorization,
$\operatorname{vec}(\mathcal Y_i)=P^\top\operatorname{vec}(\mathcal X_i)$ for
$P=U_K\otimes\cdots\otimes U_1$.
CSA constrains the PCA-like projection to a Kronecker structure; flattening
itself is lossless. The eigenupdate connects to PCA: for centered inputs,
maximizing retained energy maximizes variance; otherwise it uses second moments
\citep[Sec.~III-B]{xu2008csa}. For recognition, vectorize the core, optionally apply Fisher
LDA, and classify by 1-NN \citep[Sec.~IV-C]{xu2008csa}.
LDA uses labels to favor between-class over within-class scatter,
while 1-NN assigns the closest gallery label. These connect CSA's recognition
pipeline to supervised dimensionality reduction and instance-based learning
covered in the course.
CSA itself is unsupervised, so improved reconstruction need not improve recognition.
